\documentclass[11pt]{article}
\usepackage{amsmath,amssymb,amsthm}
\usepackage{geometry}
\usepackage{booktabs}
\usepackage{graphicx}
\usepackage{hyperref}
\usepackage{enumerate}
\usepackage{listings}
\usepackage{xcolor}

\geometry{letterpaper, margin=1in}

\lstset{
  basicstyle=\ttfamily\small,
  keywordstyle=\color{blue},
  commentstyle=\color{gray},
  stringstyle=\color{red},
  breaklines=true,
  columns=fullflexible,
  keepspaces=true,
  language=C++
}

\title{\textbf{Statistical Interpretation Layer} \\
\Large Implementation Status and Gap Analysis}
\author{Formation Engine Development Team}
\date{Version 1.0.0 --- June 28, 2025}

\begin{document}
\maketitle
\thispagestyle{empty}

\vfill
\begin{center}
\large\textit{VSEPR Simulation Program --- Internal Report} \\[0.5em]
\normalsize Derived from Section~7 audit of the statistical interpretation subsystem. \\
Documents production-ready components, conceptual-only frameworks, \\
and concrete integration tasks required for statistical rigor.
\end{center}
\vfill

\newpage
\tableofcontents
\newpage

%% ---------------------------------------------------------------
\section{Executive Summary}

An audit of the statistical interpretation layer (Section~7) reveals a split
maturity profile:

\begin{itemize}
\item \textbf{Production-ready:} Kabsch structural alignment, RMSD computation,
	  animated alignment with linear interpolation.
\item \textbf{Conceptual only:} Welford online statistics, stationarity-gated
	  convergence detection, observable tracking in integrators.
\end{itemize}

This report catalogs the current state, identifies the gaps, and proposes
prioritized integration work.

%% ---------------------------------------------------------------
\section{Production-Ready Components}

\subsection{Kabsch Algorithm --- \texttt{atomistic/core/alignment.cpp}}

The Kabsch algorithm is fully implemented and finds the optimal rigid-body
rotation $\mathbf{R}^*$ that minimizes RMSD between two $N$-atom structures:

\begin{equation}
\mathbf{R}^* = \mathbf{V}\mathbf{U}^T, \qquad
\mathbf{H} = \mathbf{U}\boldsymbol{\Sigma}\mathbf{V}^T = \sum_{i=1}^{N} \mathbf{x}_i \otimes \mathbf{r}_i
\end{equation}

SVD is computed via Jacobi iteration. The implementation handles:
\begin{itemize}
\item COM centering of both structures before alignment
\item Reflection correction ($\det(\mathbf{V}\mathbf{U}^T) < 0$)
\item Pre- and post-alignment RMSD reporting
\item Identity-rotation fallback for degenerate input ($N < 2$)
\end{itemize}

\subsection{RMSD Computation}

Standard textbook definition:
\begin{equation}
\text{RMSD} = \sqrt{\frac{1}{N}\sum_{i=1}^{N} |\mathbf{R}\mathbf{x}_i - \mathbf{r}_i|^2}
\end{equation}

Validated for conformer clustering workflows.

\subsection{Animated Alignment}

Functional with linear matrix interpolation between initial and aligned orientations.

\textbf{Known limitation:} Intermediate matrices are not proper rotations for
large angular displacements. A SLERP (quaternion) upgrade is documented but
not yet implemented.

%% ---------------------------------------------------------------
\section{Conceptual-Only Frameworks}

\subsection{Welford's Online Statistics}

Welford's one-pass recurrence avoids catastrophic cancellation when accumulating
variance over energy trajectories with large absolute values:

\begin{align}
\bar{x}_n &= \bar{x}_{n-1} + \frac{x_n - \bar{x}_{n-1}}{n} \\
M_{2,n}   &= M_{2,n-1} + (x_n - \bar{x}_{n-1})(x_n - \bar{x}_n) \\
\text{Var}(x) &= \frac{M_2}{n-1}
\end{align}

A Python reference class exists in \texttt{tools/scoring.py}. No C++
\texttt{OnlineStats} struct is present in the integrator codebase.

\subsection{Stationarity Gate / Convergence Detection}

The conceptual \texttt{ObservableTracker} monitors energy terms and temperature,
declaring convergence after $k$ consecutive samples remain within $3\sigma$ of
the running mean:

\begin{equation}
|x_{\text{new}} - \bar{x}| < 3\,\sigma(\bar{x}) + \varepsilon
\end{equation}

Currently requires manual trajectory inspection. No auto-stop mechanism exists
in Velocity Verlet, FIRE, or Langevin integrators.

%% ---------------------------------------------------------------
\section{Gap Analysis}

\begin{table}[h]
\centering
\begin{tabular}{@{}lccl@{}}
\toprule
\textbf{Component} & \textbf{Documented} & \textbf{Implemented} & \textbf{Blocking} \\
\midrule
Kabsch alignment             & Yes & Yes & --- \\
RMSD computation             & Yes & Yes & --- \\
Animated alignment           & Yes & Yes (linear) & SLERP upgrade \\
Welford \texttt{OnlineStats} & Yes & No  & Heat capacity (S4.6) \\
Stationarity gate            & Yes & No  & Auto-equilibration \\
Multi-reference Kabsch       & Yes & No  & Ensemble averaging \\
Partial (weighted) Kabsch    & Yes & No  & Backbone-only align \\
\bottomrule
\end{tabular}
\caption{Implementation status of statistical interpretation components.}
\end{table}

%% ---------------------------------------------------------------
\section{Recommended Integration Order}

\begin{enumerate}
\item \textbf{C++ \texttt{OnlineStats} struct} --- Direct port of the Python
	  reference. Required by all downstream statistical work.
\item \textbf{Integrator instrumentation} --- Add \texttt{OnlineStats} members
	  to Velocity Verlet and FIRE state structs; accumulate energy, temperature,
	  and force-magnitude variance each step.
\item \textbf{Stationarity gate in Langevin dynamics} --- Wire
	  \texttt{ObservableTracker} into the Langevin loop with configurable
	  pass count and auto-stop.
	  interpolation with quaternion geodesic interpolation on SO(3).
\item \textbf{Weighted / partial Kabsch} --- Extend alignment to support atom
	  subsets and iterative ensemble refinement.
\end{enumerate}


%% ---------------------------------------------------------------
\section{New Features: Auto-Spin and Auto-Capture}

Two new features have been added to the 3D visualization module
(\texttt{src/vis/}) to support hands-free molecular review:

\subsection{Auto-Spin}

Continuous camera orbit around the molecule center at a configurable speed.
Implemented in \texttt{AutoPilot::tick()} which calls \texttt{Camera::orbit()}
scaled by frame delta time each frame.

\begin{itemize}
\item Toggle: \textbf{S} key or ImGui checkbox
\item Speed: adjustable 10--600 px-equivalent/sec (default 120)
\item Wobble: optional vertical sinusoidal oscillation for presentation views
\end{itemize}

\subsection{Auto-Capture}

Periodic framebuffer screenshot export to PNG via \texttt{stb\_image\_write}.
Reads the GL framebuffer with \texttt{glReadPixels}, flips vertically, and
writes sequentially numbered files to a configurable output directory.

\begin{itemize}
\item Toggle: \textbf{P} key or ImGui checkbox
\item Interval: adjustable 0.25--10 seconds (default 2\,s)
\item Output: \texttt{captures/frame\_0001.png}, \texttt{frame\_0002.png}, \ldots
\item Manual trigger: `Capture Now'' button in the Auto-Pilot panel
\end{itemize}

\subsection{Integration Points}

\begin{itemize}
\item \texttt{src/vis/auto\_pilot.hpp / .cpp} --- self-contained \texttt{AutoPilot} class
\item \texttt{src/vis/window.cpp} --- ticked in all three run-loop variants
\item \texttt{src/vis/ui\_panels.cpp} --- dedicated ImGui panel with real-time controls
\end{itemize}

%% ---------------------------------------------------------------
\section{Guiding Principle}

\begin{quote}
\textit{If you cannot compute the variance of an observable, you cannot claim it
has converged. If you cannot align two structures optimally, you cannot claim
they are different. Statistical rigor is not optional---it is the foundation of
scientific validity.}
\end{quote}


\section{Offline SVG Prerendering for Reports}

The glass module now includes a complete offline report-rendering pipeline
that converts 3D molecule layouts into publication-quality SVG without
requiring an OpenGL context.

\subsection{Pipeline Summary}

\begin{enumerate}
  \item \textbf{Topology} --- \texttt{build\_topology()} constructs a
        lightweight \texttt{GlassMolecule} adjacency graph from atomic
        numbers and bond pairs.
  \item \textbf{Ring Detection} --- \texttt{detect\_rings()} identifies
        planar ring systems for correct polygon pinning.
  \item \textbf{3D Layout} --- \texttt{TopologyPrerender3D::build\_layout()}
        places atoms via BFS with ring pre-placement, chain smoothing, and
        sinusoidal arc injection.
  \item \textbf{Instance Buffers} --- \texttt{MoleculePrerenderBuilder::build()}
        converts positions into typed atom/bond instances with style flags.
  \item \textbf{SVG Export} --- \texttt{ReportRenderer::render\_svg()} performs
        software orthographic/perspective projection, depth-sorted painting,
        CPK colouring, and multi-bond rendering into a self-contained SVG string.
\end{enumerate}

\subsection{Key Design Decisions}

\begin{itemize}
  \item \textbf{No GPU dependency} --- The entire report path runs in pure C++
        with no OpenGL, GLFW, or shader requirements.
  \item \textbf{CPK colour palette} --- A 54-element Corey--Pauling--Koltun
        table maps atomic number to standard colours (H\,=\,white, C\,=\,grey,
        N\,=\,blue, O\,=\,red, etc.).
  \item \textbf{Depth shading} --- Atoms further from the camera are
        progressively dimmed, giving a pseudo-3D appearance in 2D output.
  \item \textbf{Multi-bond rendering} --- Double and triple bonds are drawn as
        parallel strokes with configurable lateral offsets.
  \item \textbf{Deterministic output} --- The layout engine uses a seeded
        xorshift RNG, so identical input always produces identical SVG.
\end{itemize}

\subsection{Usage in \LaTeX{} Documents}

\begin{lstlisting}[language=C++,basicstyle=\ttfamily\small]
#include "molecule/glass.hpp"
using namespace vsepr::glass;

auto mol     = build_topology(Z_values, bond_pairs);
auto rings   = detect_rings(mol);
auto layout  = TopologyPrerender3D{}.build_layout(mol);
auto buffers = MoleculePrerenderBuilder{}.build(mol, layout, rings);

ReportSettings rs;
rs.label_atoms = true;
ReportRenderer{rs}.write_svg("water.svg", buffers);
\end{lstlisting}

The resulting SVG can be included directly:
\begin{verbatim}
\includegraphics[width=0.5\textwidth]{water.svg}
\end{verbatim}

\\end{document}


